\IfFileExists{stacks-project.cls}{%
\documentclass{stacks-project}
}{%
\documentclass{amsart}
}

% For dealing with references we use the comment environment
\usepackage{verbatim}
\newenvironment{reference}{\comment}{\endcomment}
%\newenvironment{reference}{}{}
\newenvironment{slogan}{\comment}{\endcomment}
\newenvironment{history}{\comment}{\endcomment}

% For commutative diagrams we use Xy-pic
\usepackage[all]{xy}

% We use 2cell for 2-commutative diagrams.
\xyoption{2cell}
\UseAllTwocells

% We use multicol for the list of chapters between chapters
\usepackage{multicol}

% This is generally recommended for better output
\usepackage{lmodern}
\usepackage[T1]{fontenc}

% For cross-file-references
\usepackage{xr-hyper}

% Package for hypertext links:
\usepackage{hyperref}

% For any local file, say "hello.tex" you want to link to please
% use \externaldocument[hello-]{hello}
\externaldocument[introduction-]{introduction}
\externaldocument[conventions-]{conventions}
\externaldocument[sets-]{sets}
\externaldocument[categories-]{categories}
\externaldocument[topology-]{topology}
\externaldocument[sheaves-]{sheaves}
\externaldocument[sites-]{sites}
\externaldocument[stacks-]{stacks}
\externaldocument[fields-]{fields}
\externaldocument[algebra-]{algebra}
\externaldocument[brauer-]{brauer}
\externaldocument[homology-]{homology}
\externaldocument[derived-]{derived}
\externaldocument[simplicial-]{simplicial}
\externaldocument[more-algebra-]{more-algebra}
\externaldocument[smoothing-]{smoothing}
\externaldocument[modules-]{modules}
\externaldocument[sites-modules-]{sites-modules}
\externaldocument[injectives-]{injectives}
\externaldocument[cohomology-]{cohomology}
\externaldocument[sites-cohomology-]{sites-cohomology}
\externaldocument[dga-]{dga}
\externaldocument[dpa-]{dpa}
\externaldocument[sdga-]{sdga}
\externaldocument[hypercovering-]{hypercovering}
\externaldocument[schemes-]{schemes}
\externaldocument[constructions-]{constructions}
\externaldocument[properties-]{properties}
\externaldocument[morphisms-]{morphisms}
\externaldocument[coherent-]{coherent}
\externaldocument[divisors-]{divisors}
\externaldocument[limits-]{limits}
\externaldocument[varieties-]{varieties}
\externaldocument[topologies-]{topologies}
\externaldocument[descent-]{descent}
\externaldocument[perfect-]{perfect}
\externaldocument[more-morphisms-]{more-morphisms}
\externaldocument[flat-]{flat}
\externaldocument[groupoids-]{groupoids}
\externaldocument[more-groupoids-]{more-groupoids}
\externaldocument[etale-]{etale}
\externaldocument[chow-]{chow}
\externaldocument[intersection-]{intersection}
\externaldocument[pic-]{pic}
\externaldocument[weil-]{weil}
\externaldocument[adequate-]{adequate}
\externaldocument[dualizing-]{dualizing}
\externaldocument[duality-]{duality}
\externaldocument[discriminant-]{discriminant}
\externaldocument[derham-]{derham}
\externaldocument[local-cohomology-]{local-cohomology}
\externaldocument[algebraization-]{algebraization}
\externaldocument[curves-]{curves}
\externaldocument[resolve-]{resolve}
\externaldocument[models-]{models}
\externaldocument[functors-]{functors}
\externaldocument[equiv-]{equiv}
\externaldocument[pione-]{pione}
\externaldocument[etale-cohomology-]{etale-cohomology}
\externaldocument[proetale-]{proetale}
\externaldocument[relative-cycles-]{relative-cycles}
\externaldocument[more-etale-]{more-etale}
\externaldocument[trace-]{trace}
\externaldocument[crystalline-]{crystalline}
\externaldocument[spaces-]{spaces}
\externaldocument[spaces-properties-]{spaces-properties}
\externaldocument[spaces-morphisms-]{spaces-morphisms}
\externaldocument[decent-spaces-]{decent-spaces}
\externaldocument[spaces-cohomology-]{spaces-cohomology}
\externaldocument[spaces-limits-]{spaces-limits}
\externaldocument[spaces-divisors-]{spaces-divisors}
\externaldocument[spaces-over-fields-]{spaces-over-fields}
\externaldocument[spaces-topologies-]{spaces-topologies}
\externaldocument[spaces-descent-]{spaces-descent}
\externaldocument[spaces-perfect-]{spaces-perfect}
\externaldocument[spaces-more-morphisms-]{spaces-more-morphisms}
\externaldocument[spaces-flat-]{spaces-flat}
\externaldocument[spaces-groupoids-]{spaces-groupoids}
\externaldocument[spaces-more-groupoids-]{spaces-more-groupoids}
\externaldocument[bootstrap-]{bootstrap}
\externaldocument[spaces-pushouts-]{spaces-pushouts}
\externaldocument[spaces-chow-]{spaces-chow}
\externaldocument[groupoids-quotients-]{groupoids-quotients}
\externaldocument[spaces-more-cohomology-]{spaces-more-cohomology}
\externaldocument[spaces-simplicial-]{spaces-simplicial}
\externaldocument[spaces-duality-]{spaces-duality}
\externaldocument[formal-spaces-]{formal-spaces}
\externaldocument[restricted-]{restricted}
\externaldocument[spaces-resolve-]{spaces-resolve}
\externaldocument[formal-defos-]{formal-defos}
\externaldocument[defos-]{defos}
\externaldocument[cotangent-]{cotangent}
\externaldocument[examples-defos-]{examples-defos}
\externaldocument[algebraic-]{algebraic}
\externaldocument[examples-stacks-]{examples-stacks}
\externaldocument[stacks-sheaves-]{stacks-sheaves}
\externaldocument[criteria-]{criteria}
\externaldocument[artin-]{artin}
\externaldocument[quot-]{quot}
\externaldocument[stacks-properties-]{stacks-properties}
\externaldocument[stacks-morphisms-]{stacks-morphisms}
\externaldocument[stacks-limits-]{stacks-limits}
\externaldocument[stacks-cohomology-]{stacks-cohomology}
\externaldocument[stacks-perfect-]{stacks-perfect}
\externaldocument[stacks-introduction-]{stacks-introduction}
\externaldocument[stacks-more-morphisms-]{stacks-more-morphisms}
\externaldocument[stacks-geometry-]{stacks-geometry}
\externaldocument[moduli-]{moduli}
\externaldocument[moduli-curves-]{moduli-curves}
\externaldocument[examples-]{examples}
\externaldocument[exercises-]{exercises}
\externaldocument[guide-]{guide}
\externaldocument[desirables-]{desirables}
\externaldocument[coding-]{coding}
\externaldocument[obsolete-]{obsolete}
\externaldocument[fdl-]{fdl}
\externaldocument[index-]{index}

% Theorem environments.
%
\theoremstyle{plain}
\newtheorem{theorem}[subsection]{Theorem}
\newtheorem{proposition}[subsection]{Proposition}
\newtheorem{lemma}[subsection]{Lemma}

\theoremstyle{definition}
\newtheorem{definition}[subsection]{Definition}
\newtheorem{example}[subsection]{Example}
\newtheorem{exercise}[subsection]{Exercise}
\newtheorem{situation}[subsection]{Situation}

\theoremstyle{remark}
\newtheorem{remark}[subsection]{Remark}
\newtheorem{remarks}[subsection]{Remarks}

\numberwithin{equation}{subsection}

% Macros
%
\def\lim{\mathop{\mathrm{lim}}\nolimits}
\def\colim{\mathop{\mathrm{colim}}\nolimits}
\def\Spec{\mathop{\mathrm{Spec}}}
\def\Hom{\mathop{\mathrm{Hom}}\nolimits}
\def\Ext{\mathop{\mathrm{Ext}}\nolimits}
\def\SheafHom{\mathop{\mathcal{H}\!\mathit{om}}\nolimits}
\def\SheafExt{\mathop{\mathcal{E}\!\mathit{xt}}\nolimits}
\def\Sch{\mathit{Sch}}
\def\Mor{\mathop{\mathrm{Mor}}\nolimits}
\def\Ob{\mathop{\mathrm{Ob}}\nolimits}
\def\Sh{\mathop{\mathit{Sh}}\nolimits}
\def\NL{\mathop{N\!L}\nolimits}
\def\CH{\mathop{\mathrm{CH}}\nolimits}
\def\proetale{{pro\text{-}\acute{e}tale}}
\def\etale{{\acute{e}tale}}
\def\QCoh{\mathit{QCoh}}
\def\Ker{\mathop{\mathrm{Ker}}}
\def\Im{\mathop{\mathrm{Im}}}
\def\Coker{\mathop{\mathrm{Coker}}}
\def\Coim{\mathop{\mathrm{Coim}}}

% Boxtimes
%
\DeclareMathSymbol{\boxtimes}{\mathbin}{AMSa}{"02}

%
% Macros for moduli stacks/spaces
%
\def\QCohstack{\mathcal{QC}\!\mathit{oh}}
\def\Cohstack{\mathcal{C}\!\mathit{oh}}
\def\Spacesstack{\mathcal{S}\!\mathit{paces}}
\def\Quotfunctor{\mathrm{Quot}}
\def\Hilbfunctor{\mathrm{Hilb}}
\def\Curvesstack{\mathcal{C}\!\mathit{urves}}
\def\Polarizedstack{\mathcal{P}\!\mathit{olarized}}
\def\Complexesstack{\mathcal{C}\!\mathit{omplexes}}
% \Pic is the operator that assigns to X its picard group, usage \Pic(X)
% \Picardstack_{X/B} denotes the Picard stack of X over B
% \Picardfunctor_{X/B} denotes the Picard functor of X over B
\def\Pic{\mathop{\mathrm{Pic}}\nolimits}
\def\Picardstack{\mathcal{P}\!\mathit{ic}}
\def\Picardfunctor{\mathrm{Pic}}
\def\Deformationcategory{\mathcal{D}\!\mathit{ef}}
\usepackage{cleveref}
\crefname{originalreference}{original reference}{original references}
\Crefname{originalreference}{Original reference}{Original references}
\makeatletter
\newlabel{lemma-bbd-glueing}{{0DCB}{source}{}{sourcepointer-1}{}}
\newlabel{lemma-bbd-glueing@cref}{{[originalreference][0][]0DCB}{[1][0][]source}{}{}{}}
\makeatother
\crefname{originalreference}{original reference}{original references}
\Crefname{originalreference}{Original reference}{Original references}
\makeatletter
\newlabel{situation-locally-given}{{0D67}{source}{}{sourcepointer-2}{}}
\newlabel{situation-locally-given@cref}{{[originalreference][0][]0D67}{[1][0][]source}{}{}{}}
\makeatother
\title{Cech glueing of derived-category objects along a covering under negative-Ext and cohomological-bound vanishing}
\author{The Stacks Project Authors}
\date{stacks-project:spaces-simplicial.tex:local-snapshot-20261012; GNU FDL 1.2 or later}
\begin{document}
\maketitle
\noindent\textit{Editorial note.} All delivered mathematical text is the author's own, copied from the pinned Stacks Project chapter file sources/stacks-spaces-simplicial/pinned-original-source.tex; no statement or proof step was rewritten, reordered or removed.\par\medskip
\noindent\textit{Editorial note.} The counted proof is the author's complete argument: the '(details omitted)' at source line 5796 covers the cocycle condition for the family rho\_a: pulling back the cocycle identities of the rho\_\{ij\} (condition (1)) along the induced maps to the triple fibre products U\_i x\_X U\_j x\_X U\_k - an auxiliary pullback bookkeeping with all data defined in-span; the author's alternative 'second proof' (source lines 5830-5859), a sketch via the Cech hypercovering that ends with an explicitly omitted compatibility verification, lies outside the counted span and is not part of this card; the first sentence of the counted proof ('We urge the reader to look at the second proof in stead') is the author's framing, not an omission.\par\medskip
\noindent\textit{Editorial note.} The label lemma-bbd-glueing-cech names the counted statement, and the proof environment immediately following it in the source is delivered verbatim as its proof.\par\medskip
\noindent\textit{Editorial note.} The counted statement is carried by the author's own lemma environment in the chapter source.\par\medskip
\noindent\textit{Allowed standard prerequisites (editorial).}\begin{itemize}
\item The general glueing theorem for locally given objects. Same chapter, \textbackslash{}label\{lemma-bbd-glueing\}. Lines 5807-5828: applied with D = C and f the identity morphism; its assumptions (1), (2)(a)-(c), (3), (4) are verified.
\item The locally-given situation for a covering. Same chapter, Situation \textbackslash{}label\{situation-locally-given\}. Lines 5802-5805: the data (B, E\_U, rho\_a) are an instance of it.
\item Derived pullback along morphisms of ringed sites. Standard functoriality of derived pullback a\textasciicircum{}*. Lines 5788-5800: E\_U = a\_U\textasciicircum{}* E\_\{i(U)\} and the isomorphisms rho\_a from the rho\_\{ij\}.
\item Vanishing of negative self-Exts descends along restriction. Functoriality of SheafExt under derived pullback. Lines 5819-5821: negative Ext vanishing for E\_i implies the same for every restriction, feeding assumption (3).
\end{itemize}\medskip
\ifdefined\hypertarget\hypertarget{sourcepointer-1}{}\fi
\noindent\textit{Source reference.} Original reference 0DCB: sources/stacks-spaces-simplicial/pinned-original-source.tex, source lines 5630--5630. Prerequisite pointer invoked at proof line(s) 5776; it names a standard fact of the same chapter that the counted argument uses, and every step of the counted proof built on it is retained verbatim in the delivered body.\par\medskip
\ifdefined\hypertarget\hypertarget{sourcepointer-2}{}\fi
\noindent\textit{Source reference.} Original reference 0D67: sources/stacks-spaces-simplicial/pinned-original-source.tex, source lines 5552--5552. Prerequisite pointer invoked at proof line(s) 5801; it names a standard fact of the same chapter that the counted argument uses, and every step of the counted proof built on it is retained verbatim in the delivered body.\par\medskip
\begin{lemma}
\label{lemma-bbd-glueing-cech}
\begin{reference}
Email of Martin Olsson dated Sep 9, 2021.
\end{reference}
Let $(\mathcal{C}, \mathcal{O}_\mathcal{C})$ be a ringed site.
Assume $\mathcal{C}$ has fibre products.
Let $\{U_i \to X\}_{i \in I}$ be a covering in $\mathcal{C}$.
For $i \in I$ let $E_i$ be an object of $D(\mathcal{O}_{U_i})$ and
for $i, j \in I$ let
$$
\rho_{ij} : E_i|_{\mathcal{C}/U_{ij}} \longrightarrow E_j|_{\mathcal{C}/U_{ij}}
$$
be an isomorphism in $D(\mathcal{O}_{U_{ij}})$ where
$U_{ij} = U_i \times_X U_j$. Assume
\begin{enumerate}
\item the $\rho_{ij}$ satisfy the cocycle condition on
$U_i \times_X U_j \times_X U_k$ for all $i, j, k \in I$,
\item $\SheafExt^p_{\mathcal{O}_{U_i}}(E_i, E_i) = 0$
for all $p < 0$ and $i \in I$, and
\item there exists a $t \in \mathbf{Z}$ such that
$H^p(E_i) = 0$ for $p < t$ and all $i \in I$.
\end{enumerate}
Then there exists a unique pair $(E, \rho_i)$ where $E$ is an object of 
$D(\mathcal{O}_X)$ and $\rho_i : E|_{U_i} \to E_i$ are isomorphisms in
$D(\mathcal{O}_{U_i})$ compatible with the $\rho_{ij}$.
\end{lemma}

\begin{proof}[First proof]
In this proof we deduce the lemma from the very general
Lemma \ref{lemma-bbd-glueing}. We urge the reader to look
at the second proof in stead.

\medskip\noindent
We may replace $\mathcal{C}$ with $\mathcal{C}/X$.
Thus we may and do assume $X$ is the final object of $\mathcal{C}$
and that $\mathcal{C}$ has all finite limits.

\medskip\noindent
Let $\mathcal{B}$ be the full subcategory of $\mathcal{C}$
consisting of $U \in \Ob(\mathcal{C})$ such that there exists
an $i(U) \in I$ and a morphism $a_U : U \to U_{i(U)}$.
We denote $E_U = a_U^*E_{i(U)}$ in $D(\mathcal{O}_U)$
the pullback (restriction) of $E_i$ via $a_U$.
Given a morphism $a : U \to U'$ of $\mathcal{B}$
we obtain a morphism
$(a_{U'} \circ a, a_U) : U \to U_{i(U')} \times_X U_{i(U)} = U_{i(U')i(U)}$
and hence an isomorphism
$$
\rho_a :
a^*E_{U'} = a^*a_{U'}^*E_{i(U')}
\xrightarrow{(a_{U'} \circ a, a_U)^*\rho_{i(U')i(U)}}
a_{U}^*E_{i(U)} = E_{U}
$$
in $D(\mathcal{O}_U)$. The data $\mathcal{B}, E_U, \rho_a$
are as in Situation \ref{situation-locally-given}; the
isomorphisms $\rho_a$ satisfy the cocycle condition
exactly because of condition (1) in the statement of the lemma
(details omitted).

\medskip\noindent
We are going to apply Lemma \ref{lemma-bbd-glueing} with $\mathcal{B}$,
$E_U$, $\rho_a$ as above and with $\mathcal{D} = \mathcal{C}$ and
$f : \mathcal{C} \to \mathcal{D}$ the identity morphism.
Assumptions (1) and (2)(a) of Lemma \ref{lemma-bbd-glueing} we have seen above.
Assumption (2)(b) of Lemma \ref{lemma-bbd-glueing} is clear.
Assumption (2)(c) of Lemma \ref{lemma-bbd-glueing} holds because
$\{U_i \to X\}$ is a covering\footnote{In fact, it would suffice
if the map $\coprod_{i \in I} h_{U_i} \to h_X$ becomes
surjective on sheafification and the lemma holds in this case
with the same proof.}.
Assumption (3) of Lemma \ref{lemma-bbd-glueing} holds
because we have assumed the vanishing of all negative Ext sheaves
of $E_i$ which certainly implies that for any object $U$ lying over $U_i$
the negative self-Exts of $E_i|_U$ are zero.
Assumption (4) of Lemma \ref{lemma-bbd-glueing} holds because
we have assumed the cohomology sheaves of each $E_i$ are zero
to the left of $t$.

\medskip\noindent
We obtain a unique solution $(E, \rho_U)$.
Setting $\rho_i = \rho_{U_i}$ the lemma follows.
\end{proof}

\end{document}
